\section{Proofs for reference-centered risk qualification}\label{supp:centeredrisk92}
\subsection{Centering and the original Bellman optimum}
For any bounded $z$, substituting $f_{t+1}+z$ in \eqref{eq:dualrisk92} and subtracting the finite number $\rho(f_{t+1})$ proves \eqref{eq:exact-center92}. Each dual objective at $f_{t+1}$ is at most its supremum, giving $k^f\ge0$. Approximating that supremum gives $\inf_q k^f(q)=0$ without an attainment assumption. Since each $q$ is a probability, adding a constant $b$ to $z$ adds $b$ to $\mathcal E^f(z)$. If $z\le w$, every dual objective at $z$ is no larger than its counterpart at $w$, proving monotonicity. The inequalities $w-\|z-w\|_\infty\le z\le w+\|z-w\|_\infty$ prove nonexpansiveness.

\begin{proof}[Proof of Theorem~\ref{thm:centeredrisk92}]
Let $d_t^\pi=V_t^\pi-f_t$ and $d_t^*=V_t^*-f_t$. The original recursions and \eqref{eq:exact-center92} give the exact identities
\begin{align}\label{eq:exact-error-rec92}
 d_t^\pi(x)&=r_t(x,\pi_t(x))+
            \beta_t\mathcal E_{t,x,\pi_t(x)}^f(d_{t+1}^\pi),\notag\\
 d_t^*(x)&=\inf_{a\in A_t(x)}
            \{r_t(x,a)+\beta_t\mathcal E_{t,x,a}^f(d_{t+1}^*)\}.
\end{align}
Both have terminal value $g-f_T$. Suppose $d_{t+1}^\pi\le U_{t+1}$. Monotonicity, nonnegative $\beta_t$, the upper residual inequality and \eqref{eq:potential-cert92} imply $d_t^\pi\le U_t$. Similarly $L_{t+1}\le d_{t+1}^*$ implies $L_t\le d_t^*$. Backward induction establishes both statements. Feasibility of the supplied complete policy implies $V_t^*\le V_t^\pi$, giving \eqref{eq:potential-loss92}. The argument takes place at each fixed parameter vector; applying it to every vector satisfying the premises proves the class statement without assembling incompatible continuations.

Exact equality recursions coincide with \eqref{eq:exact-error-rec92}, proving the exactness assertion by induction. For the scalar comparison, set $a_t=\ell_t+b_t^-$ and $b_t=u_t+\gamma_t+b_t^{\pi,+}$ from the signed assumptions. Every original action has residual at least $a_t$, since its numerical action value is at least the numerical Bellman infimum. The implemented action has residual at most $b_t$. Define $\underline h_T=d^-$, $\overline h_T=d^+$, and
\[
 \underline h_t=a_t+\beta_t\underline h_{t+1},\qquad
 \overline h_t=b_t+\beta_t\overline h_{t+1}.
\]
Since $\mathcal E^f(b)=b$ for constants, these are feasible lower and upper certificates. Exact backward evaluation is order preserving, so $L_t\ge\underline h_t$ and $U_t\le\overline h_t$ for exact residuals and terminal discrepancies. Their date-zero difference is at most \eqref{eq:signed-loss91}.

For two policy supersolutions, their pointwise minimum remains a supersolution: the operator at the minimum of their successor bounds is at most its value at either bound. The corresponding pointwise maximum of two optimum subsolutions remains a subsolution by the reverse argument. Finite combinations therefore preserve an already valid scalar guarantee. No claim about improved bounds from an uncontrolled numerical optimizer is needed.
\end{proof}

\subsection{Risk occupations and certificate orientation}
\begin{proof}[Proof of Proposition~\ref{prop:riskflow92}]
For fixed $\pi$, \eqref{eq:exact-error-rec92} is the Bellman equation for a finite-horizon maximizer selecting a risk index $j$ after observing the current state. Its reward is $r_{ts}-\beta_tk_{tsj}$, transition vector is $q_{tsj}$, and terminal reward is $g-f_T$. Finite maxima have measurable maximizers. Backward induction establishes a Markov maximizing choice at each date and state.

Any nonanticipating choice rule, including a history-dependent one, induces occupation masses satisfying \eqref{eq:riskflow92}. Conversely, for a feasible flow, choose $j$ at $(t,s)$ with probabilities $y_{tsj}/m_{ts}$ when $m_{ts}>0$; at zero-mass states choose an arbitrary feasible index. Induction on $t$ recovers the declared masses. The corresponding expected sum of rewards is exactly the objective in \eqref{eq:riskflow-value92}. Thus the flow maximum equals the centered fixed-policy value. This representation is under the adversarial risk kernels and must not be called the physical occupation law.

Multiply each dual inequality by $B_ty_{tsj}$ and sum. The flow equalities cancel intermediate potential terms, leaving the primal objective at most $\mu^{\mathsf T}U_0$. The argument allows a zero discount: it uses $B_{t+1}=B_t\beta_t$ but never divides by either factor. Feasible primal and dual witnesses with equal objectives certify the maximum. For an optimum subsolution, a separate feasible risk witness for every original action supplies a lower bound on that action's supremum; a minimization over a restricted sampled action set does not supply the required all-action inequality.
\end{proof}

\subsection{Kernel exclusion and its sharp boundary}
For any $z\in\mathcal D$, the difference between the affine pieces indexed by $q$ and $q_0$ is
\[
 (q-q_0)^{\mathsf T}z-k^f(q)+k^f(q_0).
\]
A strict negative upper bound excludes $q$ uniformly. This proves \eqref{eq:kernel-cut92}. Maximizing coordinate by coordinate gives the box formula. For a span ball, the coefficients sum to zero; adding a constant to $z$ changes nothing. Normalize $\min_i z_i=0$ and assign $s$ to the positive coefficients and zero to the negative ones. The maximum is $s\sum_i(q_i-q_{0i})_+=s\|q-q_0\|_1/2$.

As a sharp example, take two equally likely outcomes and average value at risk at level $1/2$. Then $\rho(v)=\max(v_1,v_2)$ and the dual vertices are $e_1,e_2$. At $f=(0,1)$ their centered penalties are $1,0$. For $z=(2,0)$, the centered value is $1$. Keeping only the maximizing kernel $e_2$ from the reference gives $0$, which is invalid as an upper bound; dropping penalties gives $2$, which is conservative but not exact. On the span ball of radius $1/2$, the first kernel is safely excluded. At radius $1$, $z=(1,0)$ makes the two pieces equal, so the strict exclusion must retain the boundary. This example proves both the necessity of allowing the risk kernel to change and sharpness of the strict threshold convention.

\subsection{Entropic duality without changing the objective}
Let $p$ be a probability measure, $f,z$ bounded, and $\lambda>0$. Boundedness makes $Z_f=\int e^{\lambda f}dp$ strictly positive and finite, and makes $p_f$ equivalent to $p$ on its support. For $q\ll p$ with finite relative entropy,
\begin{align*}
 \lambda^{-1}{\rm KL}(q\Vert p_f)
 &=\lambda^{-1}{\rm KL}(q\Vert p)-\int f\,dq
                                      +\lambda^{-1}\log Z_f.
\end{align*}
If $q$ places mass outside the support of $p$, its penalty is infinite. For any bounded $v$, define $p_v=e^{\lambda v}p/Z_v$. Nonnegativity of ${\rm KL}(q\Vert p_v)$ gives
\[
 \int v\,dq-\lambda^{-1}{\rm KL}(q\Vert p)
 \le\lambda^{-1}\log Z_v,
\]
with equality at $q=p_v$. This proves the variational representation, its centered penalty, and \eqref{eq:entropic-center92}. It also shows exactly why using the fixed tilt $p_f$ as an expectation is not generally exact: the centered evaluation remains entropic under $p_f$, rather than linear. For $\lambda\downarrow0$, bounded convergence gives the expectation limit. Exponential and logarithmic evaluations in a numerical certificate require outward arithmetic; the rational average-value-at-risk experiments are not evidence of exact floating-point entropic computation.

\subsection{Occupation-sensitive precision and a strict improvement}
For fixed penalties and transitions, write the flow objective as $A(y)$ and the additional nonnegative residual contribution as $E(y)=\sum_{t,s}B_te_{ts}m_{ts}$. If $H=\max_y A(y)$ and $H_e=\max_y\{A(y)+E(y)\}$, then
\[
 0\le H_e-H\le\max_y E(y).
\]
The upper inequality follows from $A(y)\le H$ for every feasible $y$; the lower follows from $E(y)\ge0$. Since each occupation layer has total mass one, $\max_y E(y)\le\sum_t B_t\max_s e_{ts}$. This establishes \eqref{eq:occupation-budget92}. The argument assumes the penalties are held fixed while residual uncertainty is enlarged. Altering the underlying stored continuation changes the penalties and requires requalification.

For a strict horizon example, let two states be absorbing under both actions, let action zero cost zero and action one cost one in both states, let $g=0$, and let all discounts equal one. Every normalized monetary risk evaluation of these deterministic transitions agrees with evaluation at the absorbing state. Put $f_t(s)=(T-t)s$ for $s\in\{0,1\}$ and choose action zero everywhere. Each date's implemented and minimum residuals are $(0,-1)$; the signed scalar account charges one at each date, yielding $T$. The exact state-dependent lower and upper recursions both equal $-(T-t)s$, giving zero policy loss at every state. There are two feasible actions and every action is checked. No state-dependent bias was cancelled without propagation: its entire state dependence was carried through the original transition at every date. This is consistent with, and sharpens rather than contradicts, the statewise-bias warning in Section~\ref{sec:risk91}.

Finally, the finite backward recursion evaluates each stored transition support once for each retained risk piece at each date and state-action pair. Summing these support sizes yields the stated arithmetic count. A direct descendant recursion repeatedly evaluates the same subproblem; removing that repetition does not establish a polynomial bound in the number of continuous states, risk pieces, acquisition cells or rational bit length. The comparison is a representation-conditioned resource statement, and the full construction and support-qualification work remains separately charged.
